\documentclass{situationist-preprint}

\title{Global Justice and Institutional Duties}
\msid{POL-13-a}
\msfamily{POL-13}
\msdiscipline{Political Philosophy}
\mstype{Principal manuscript}
\mskeywords{global justice; beneficence; statism; international institutions}

\begin{document}
\maketitle

\begin{abstract}
What do the affluent owe the global poor? ETH-06 argues that affluent agents must give a fixed share of their income, ten per cent, to the collective task of meeting basic need, a share set so that need is met even if only those holding a tenth of affluent income comply. We argue that this duty applies globally, and that its fulfilment requires international institutions. The duty is grounded in need and in the capacity to meet it, neither of which stops at borders. But a duty to contribute a share to a collective task requires assessment, allocation and assurance that only institutions supply. Because the share is robust to partial compliance, the institutions need only secure compliance among those holding one tenth of affluent income, not universal compliance.
\end{abstract}

\section{The problem}

\citet{Singer1972} argued that distance makes no difference to the duty to aid; \citet{Beitz1979} and \citet{Pogge2002} extended distributive justice to the world. Statists deny the extension: \citet{Nagel2005} argued that the demands of justice arise only among those who share a coercive political order, and \citet{Rawls1999} limited international obligation to a duty of assistance. We take ETH-06 as given and ask whether its duty holds between compatriots only, and what its fulfilment requires.

\section{Framework}

\begin{assumption}[The duty, after ETH-06]\label{a:duty}
Each affluent agent has a duty to contribute a fixed share of income, ten per cent, to the collective task of meeting basic need; the share is set so that basic need is met if the complying agents hold at least one tenth of the income of the affluent.
\end{assumption}

\begin{definition}\label{d:inst}
An institution \emph{serves} the collective task if it (i) assesses where basic need lies and what meeting it costs, (ii) allocates contributions so that they meet need without duplication or gaps, and (iii) assures each contributor that her share is part of a scheme that reaches the compliance floor.
\end{definition}

\section{The argument}

\begin{proposition}[Globality]\label{p:global}
The duty of Assumption~\ref{a:duty} is owed to all in basic need, wherever they live.
\end{proposition}

\begin{proof}
The duty is grounded in two facts: that some are in basic need, and that the affluent can meet it at moderate cost. Neither depends on a shared coercive order. Nagel's argument concerns egalitarian justice; he accepts a humanitarian duty to relieve extreme deprivation that holds between all persons \citep{Nagel2005}, and the duty of ETH-06 is of this kind.
\end{proof}

\begin{proposition}[Institutional fulfilment]\label{p:inst}
The duty of Assumption~\ref{a:duty}, applied globally, can be fulfilled only through institutions that serve the collective task.
\end{proposition}

\begin{proof}
The duty is a duty to contribute to a collective task, not merely to transfer a sum. An individual acting alone cannot know where basic need lies or what meeting it costs, cannot coordinate her gift with millions of others, and cannot know whether the compliance floor is reached. These are conditions (i)--(iii) of Definition~\ref{d:inst}. Within a state, fiscal and welfare institutions supply them; globally, only international institutions can, whether public agencies and multilateral funds or private agencies under public standards of accountability.
\end{proof}

\begin{corollary}[The compliance floor]\label{c:floor}
Institutions that serve the global task need only secure compliance among affluent agents holding one tenth of affluent income.
\end{corollary}

\begin{proof}
By Assumption~\ref{a:duty}, need is met at that level of compliance.
\end{proof}

Corollary~\ref{c:floor} makes the requirement feasible. Institutions securing full compliance, as a full-compliance share would require \citep{Murphy2000}, would need coercive powers over the affluent of every country; institutions that need reach only a tenth can rely on voluntary participation, national aid budgets and existing agencies.

\section{An objection}

\begin{objection}[No institutions, no duty]
If fulfilment requires institutions that are absent or inadequate, the duty lapses.
\end{objection}

\begin{reply}
It changes form: the affluent must contribute through the best available agencies and support the creation of better ones, a version of the natural duty to further just arrangements not yet established \citep[§19]{Rawls1971}.
\end{reply}


\section{Conclusion}

The duty of ETH-06 is owed to the global poor, as a share in a task that only international institutions can carry out, and that they can carry out at a compliance of one in ten.

\begin{thebibliography}{99}\small
\bibitem[Beitz(1979)]{Beitz1979} Beitz, C.~R. (1979). \emph{Political Theory and International Relations}. Princeton University Press.
\bibitem[Murphy(2000)]{Murphy2000} Murphy, L.~B. (2000). \emph{Moral Demands in Nonideal Theory}. Oxford University Press.
\bibitem[Nagel(2005)]{Nagel2005} Nagel, T. (2005). The problem of global justice. \emph{Philosophy \& Public Affairs}, 33(2), 113--147.
\bibitem[Pogge(2002)]{Pogge2002} Pogge, T. (2002). \emph{World Poverty and Human Rights}. Polity Press.
\bibitem[Rawls(1971)]{Rawls1971} Rawls, J. (1971). \emph{A Theory of Justice}. Harvard University Press.
\bibitem[Rawls(1999)]{Rawls1999} Rawls, J. (1999). \emph{The Law of Peoples}. Harvard University Press.
\bibitem[Singer(1972)]{Singer1972} Singer, P. (1972). Famine, affluence, and morality. \emph{Philosophy \& Public Affairs}, 1(3), 229--243.
\end{thebibliography}

\end{document}
